\documentclass[../textbook.tex]{subfiles}
\begin{document}
\bookchapter{领域应用}{ch:domain-applications}

领域应用不是给通用模型增加一个行业名称。它需要明确要改善的工作、可依赖的知识、允许改变的业务状态，以及谁能判定结果合格。同一模型在文本归纳上表现良好，并不能推出它能独立完成合同判断、设备处置或专业决策。应用设计应从具体任务及其证据要求出发。

本章沿用\bookxref{ch:rag-systems}与\bookxref{ch:agent-runtime}的机制，讨论知识服务、业务流程、工业与专业服务。NIST AI RMF 1.0 将使用情境、测量与持续管理纳入系统生命周期\cite{nist2023airmf}；本章据此强调任务级证据，而不把框架名称作为应用已合格的证明。

\section{任务定义及方法选择}

\subsection{业务任务定义}
\term{任务单元}{Task Unit}{}是可计数、可验收的一次业务工作，例如从一份指定版本的合同抽取付款条件，或为一个告警建立内容完整的待审工单。任务单元需要固定输入边界、输出结构、容许失败、人工职责和结束条件。只写“提升效率”无法确定分母，也无法区分更快完成与更快产生待返工内容。

设输入为 $x$，外部事实状态为 $z$，系统输出或动作为 $a$，业务损失为 $L(a,z)$。理想决策可写为
\begin{equation}
 a^*(x)=\arg\min_{a\in\mathcal A(x)}\mathbb E[L(a,z)\mid x],
 \label{eq:domain-decision}
\end{equation}
其中 $\mathcal A(x)$ 是权限、业务规则和资源约束允许的动作集合。概率估计、损失函数与可行动作都需要领域依据；模型生成概率不是业务正确概率，不能直接代入式\eqref{eq:domain-decision}。本章主要符号见表\ref{tab:domain-applications-symbols}。

\begin{table}[htbp]
\centering
\caption{领域应用的主要符号与计量对象。}
\label{tab:domain-applications-symbols}
\begin{tabularx}{\linewidth}{@{}lX@{}}
\toprule
符号 & 含义\\
\midrule
$x,z,a$ & 输入、真实业务状态与输出或动作。\\
$L(a,z),\mathcal A(x)$ & 业务损失及允许动作集合。\\
$Y(1),Y(0)$ & 同一任务在辅助流程与原流程下的潜在结果。\\
$T_0,T_a,T_r,T_w$ & 原处理时间、辅助操作时间、复核时间与返工时间。\\
$p,C_e,C_h$ & 经校准的正确概率、错误损失与人工处置成本。\\
$N,n_s,n_h$ & 任务总数、合格完成数与人工干预数。\\
\bottomrule
\end{tabularx}
\end{table}

\subsection{提示、检索、微调及工具的选择}
\begin{table}[htbp]
\centering
\caption{业务缺口、候选机制及其验证边界。}
\label{tab:domain-applications-mechanisms}
\begin{tabularx}{\linewidth}{@{}p{.24\linewidth}p{.33\linewidth}X@{}}
\toprule
主要缺口 & 对应机制 & 必须验证的边界\\
\midrule
目标或格式不清楚 & 指令、示例与结构化输出 & 是否覆盖难例，解析有效是否伴随内容正确。\\
事实不在上下文或更新频繁 & 检索与证据组装 & 权限、版本、召回及逐声明支持。\\
稳定任务行为与表达需要适配 & 监督微调或参数高效微调 & 数据代表性、遗忘、域外表现及维护成本。\\
需要精确计算或实时状态 & 受控计算与业务查询工具 & 单位、对象身份、时间、新鲜度和错误传播。\\
需要改变业务状态 & 有授权的动作执行与持久运行 & 幂等、审批、回执、补偿与恢复。\\
\bottomrule
\end{tabularx}
\end{table}

表\ref{tab:domain-applications-mechanisms}中的机制可以组合。微调能改善术语与任务行为，却不能自动替代每日变化的库存读取；检索能补充文档，却不能替代数据库约束；增加动态规划也不能修复身份映射错误。方法选择首先对应已定位的缺口，再通过对照实验判断是否必要。

\section{领域知识及数据契约}

\subsection{事实的时间对象及单位}
同一句“库存为100”缺少物料、仓库、批次、单位、状态和时点，通常不足以驱动动作。领域数据应把这些条件作为字段，而非期望模型从上下文猜测。知识文档还需要适用范围、生效日期、废止关系和来源等级；历史有效与当前有效必须区别。

\term{领域本体}{Domain Ontology}{}描述实体类型及其关系；应用可以从稳定的术语映射、枚举和主数据关系开始，并不要求先建立庞大知识图谱。例如设备标识、工位、维护计划与故障代码的映射必须来自受信主数据。检索相似的设备名称不能替代精确对象解析。

表格与文档抽取要保留单位、列标题、脚注和跨页关系。金额的币种、数量的计量单位、日期的时区或业务日口径都会影响结论。计算交给具有明确输入类型与规则的工具，模型负责解释经核对的结果；工具同样可能接收错误输入，所以必须核查抽取证据及单位转换。

\subsection{专家反馈的可靠性}
专家标注应记录问题、证据版本、判断标准和分歧处理。同一文本在不同业务时点可能有不同答案；多数票也可能掩盖专业分工差异。首先澄清任务定义，再建立双人复核或仲裁，而不是机械合并所有意见。

训练集与评估集按文档、客户、设备或时间等真实关联单位分组，避免同一事实的不同切块跨集合泄漏。上线后的失败样本很有价值，但通常偏向被发现的问题；不能直接将其分布视为全部业务分布。主动收集成功、拒答、人工改写和未被采用的建议，才能理解完整行为。

\section{知识及文档密集型应用}
\optional
\subsection{企业知识服务}
知识助手应将回答拆成可以核验的声明，每项声明关联到证据的局部位置。检索到文件名并不证明具体结论得到支持。权限过滤、版本冲突和无答案路径是正常流程：若政策版本不同，应呈现适用范围与冲突，而非让模型静默混合。

可以分别测量材料召回、有效证据覆盖、声明支持、回答完整与任务完成。用户点赞同时受措辞与速度影响，不能独立作为事实准确性指标。多轮对话还应保存查询时点，避免把前轮已过时状态当作当前事实。

\subsection{合同审阅及财务分析}
合同应用可先把任务限定为条款定位、字段抽取、版本差异和待审问题清单。每个结果保留原文位置、条件与例外；风险标签是供专业人员复核的工作产物，不能仅凭模型置信措辞变成法律结论。条款的行业、地域、主体和生效条件可能改变适用性，因此应由领域流程确定判断标准。

财务分析首先区分已披露事实、计算所得指标和预测假设。抽取金额必须保留币种、期间、合并口径与单位；同比比较要确认基期和重述版本一致。设当期值为 $v_1$、基期值为 $v_0>0$，增长率 $\frac{v_1-v_0}{v_0}$ 的计算简单，错误常来自两个数不具有可比性。基期为零或符号变化时，应另行解释而非输出无意义百分比。

\begin{bookinsight}
对专业文档，应用职责可以限定为信息整理与证据定位，并在指定任务集上验证其可靠性；是否允许进一步作出专业判断，需要独立的任务验证和责任安排。本章讨论系统设计，不据此给出具体法律、投资或医疗处置意见。
\end{bookinsight}

\section{业务流程及工业应用}
\optional
\subsection{客服、工单及办公流程}
客服任务应区分回答政策、解释订单状态和执行退款等动作。回答正确但重复创建工单仍属于端到端失败；工单创建成功但缺少必要信息也可能增加后续处理成本。流程指标因此需要覆盖首次解决、重开、转交、返工和最终业务结果。

自动分派可先生成结构化候选，规则层检查队列范围、责任归属与服务级别，再由业务系统条件更新。办公流程中的审批人、金额阈值和组织权限来自主数据及规则系统，不应由模型臆测。跨步骤恢复机制见\bookxref{ch:agent-runtime}[中的持久运行系统]。

\subsection{制造业知识及 ERP/MES 集成}
\term{企业资源计划}{Enterprise Resource Planning}{ERP}与\term{制造执行系统}{Manufacturing Execution System}{MES}承载不同粒度的业务状态。大模型应用可以检索作业指导书、汇总异常记录、解释经校验的生产数据，并准备待审业务动作，但不能把自然语言建议直接等同于设备控制指令。

例如解释设备告警时，系统需要关联设备标识、固件或配置版本、告警时间、工况、维护记录和适用指导书。日志中的先后发生不自动构成因果关系，模型提出的故障假设应与已观察事实分列。若需修改工单或维护计划，使用版本条件与回执；若涉及实时控制，则必须遵循控制系统自身的确定性约束与验证边界。

\begin{figure}[htbp]
\centering
\begin{tikzpicture}[font=\small,node distance=6mm,
 box/.style={draw,rounded corners,align=center,text width=30mm,minimum height=12mm}]
\node[box] (sources) {领域事实\\主数据、文档、事件};
\node[box,right=of sources] (knowledge) {知识与计算层\\检索、规则、统计};
\node[box,right=of knowledge] (model) {模型层\\解释与动作提案};
\node[box,below=of model] (control) {执行控制层\\授权、版本、审批};
\node[box,left=of control] (business) {业务系统\\条件更新与回执};
\node[box,left=of business] (eval) {验证与反馈\\结果、返工、失败};
\draw[->] (sources)--(knowledge);\draw[->] (knowledge)--(model);
\draw[->] (model)--(control);\draw[->] (control)--node[above,font=\scriptsize]{仅允许}(business);
\draw[->] (business)--(eval);\draw[->] (eval)--(sources);
\end{tikzpicture}
\caption{领域应用以事实、受控执行和业务反馈形成闭环。}\label{fig:domain-system}
\end{figure}

图\ref{fig:domain-system}将知识解释与业务更新分开；只有经过授权和前置状态检查的动作才能产生业务回执，模型层不能越过执行控制层。

\section{专业服务及研发应用}
\optional
教育辅导的目标可以是解释概念、提供分层提示或帮助发现错误。直接给出答案与促进学习不是同一指标，应观察学生后续独立完成能力，而非仅观察当次对话满意度。评价需要区分年龄、先修知识、题目泄漏及帮助程度。

医疗信息辅助可限定为资料检索、记录整理和供专业人员核对的摘要。资料的适用人群、时点、证据等级和缺失信息必须保留；摘要遗漏与编造同样重要。系统是否可以参与具体决策，需要超出文本基准的场景验证，不能由一般问答准确率推导。

科研知识服务应区分论文报告、复现结果与作者推测，保留实验条件及不一致证据。引用存在不等于论点得到支持；自动文献综述应避免只收集支持最初假设的材料。生成新的研究假设可以有价值，但其身份始终是假设，必须经过独立验证。

代码研发助手的任务可包括定位、修改、测试与评审。编译通过只说明满足一部分语法或类型约束；测试通过也只支持测试实际覆盖的行为。应记录变更范围、依赖、真实运行证据和未覆盖场景。外部仓库内容与执行输出属于资料，不能扩张凭证或发布权限。

这些领域共享一个原则：将系统能力限定为可以证明的任务行为，再根据证据扩展。不同领域的责任与误差成本不相同，不能共用一个通用通过率作为所有上线门槛。

\section{业务效果的计算及归因}

\subsection{复核返工成本}
\phantomsection\label{q:ch36:example:01}


设原流程每任务耗时 $T_0$，辅助流程包含操作 $T_a$、复核 $T_r$ 与平均返工 $T_w$，则平均净节省为
\begin{equation}
 \Delta T=\mathbb E[T_0]-\mathbb E[T_a+T_r+T_w].
 \label{eq:domain-timesave}
\end{equation}
返工为零只适用于确实没有返工的任务，不能因为日志未记录就填零。若生成节省时间，但复核更慢，总收益可能为负。

考虑一个明确的构造算例：每任务原需12分钟，辅助操作4分钟，复核3分钟，20\%任务还需10分钟返工。新平均时间为 $4+3+0.2\times10=9$ 分钟，净节省3分钟，而非仅用 $12-4=8$ 分钟。100个任务节省300分钟；这还没有扣除维护、评估和部署投入。

\subsection{自动化覆盖及错误成本}
\phantomsection\label{q:ch36:example:02}


设某任务经独立校准的正确概率为 $p$，错误成本为 $C_e$，人工处理成本为 $C_h$。在简化假设“正确自动处理无额外成本，人工结果无错误”下，自动处理的期望损失为 $(1-p)C_e$，因此只有
\begin{equation}
 p\ge 1-\frac{C_h}{C_e}
\end{equation}
时自动处理的期望成本不高于人工。若自动处理也有成本 $C_a$、人工仍有残余风险，则必须把两者加入比较。这是阈值机制的解释，不是可直接投入业务的通用分数门槛。

提高阈值通常降低自动覆盖率。应一起报告覆盖率与被自动处理部分的错误率，并核查被转交任务是否使人工负担更重。模型自报“90\%确定”并不满足上述校准前提；需要用历史标注与独立评估建立置信估计。

\subsection{对照实验及分母}
令 $Y_i(1),Y_i(0)$ 表示同一任务在辅助和原流程下的潜在结果，目标平均效应为 $\mathbb E[Y(1)-Y(0)]$。同一任务通常只能观察一种处理结果，因此简单上线前后比较还混有任务难度、人员和季节变化。随机分配或适当的配对设计有助于隔离这些因素，统计方法见\bookxref{ch:evaluation-statistics}。

如果同一工作人员同时接触两种流程并把学到的方法用于对照组，可能发生干预扩散；如果任务来自同一客户或设备，样本也可能相关。应按实际关联层级分组或在分析中处理相关性，不能把所有消息当作独立任务扩大样本量。

\begin{table}[htbp]
\centering
\caption{领域应用评估指标及其统计口径。}
\label{tab:domain-applications-metrics}
\begin{tabularx}{\linewidth}{@{}p{.28\linewidth}p{.31\linewidth}X@{}}
\toprule
指标 & 定义 & 必须说明的口径\\
\midrule
任务成功率 & $\frac{n_s}{N}$ & 合格标准、拒答与放弃是否计入。\\
人工干预率 & $\frac{n_h}{N}$ & 一次或多次干预如何计数。\\
平均处理时长 & 总任务处理时间除以任务数 & 是否包含等待、复核和返工。\\
单位合格任务成本 & 总流程成本除以 $n_s$ & 失败消耗与人工成本是否包含。\\
自动覆盖率 & 自动结束任务数除以 $N$ & 自动结束是否同时满足质量门槛。\\
\bottomrule
\end{tabularx}
\end{table}

表\ref{tab:domain-applications-metrics}采用统一分母解释指标。例如提交1000项任务，700项自动结束，其中630项合格，另有270项经人工后合格、30项仍失败。自动覆盖率70\%，自动部分合格率90\%，整体成功率90\%。若只报告“自动完成700项”，会隐去70项自动错误。这个算例说明分母如何改变解释，并不代表某系统的测量结果。

\section{失败归因及改进}

失败可以来自对象识别、数据新鲜度、检索、解释、工具参数、执行、复核或任务定义。应保存最小必要证据以重建失败，不用“模型不够强”替代定位。一次改进尽量对应明确机制，并在冻结对照集上观察回归：更长上下文可能改善材料覆盖却增加错误混合，更强重排也可能压低少数来源的可见性。

\begin{bookalgorithm}
\textbf{领域应用的证据闭环。}
\begin{enumerate}
\item 定义任务单元、输出契约、允许动作与业务损失，指定结果判定者。
\item 建立带来源、版本、权限和时间的领域样本，并按关联对象与时间划分训练、开发和评估。
\item 保留原流程基线，选择与已定位缺口对应的提示、检索、适配或工具机制。
\item 记录端到端结果、人工干预、返工、成本及失败原因；对拒答和未知结果单独保留状态。
\item 以对照或配对方案评估收益，同时检查关键切片和业务约束，避免用总体均值掩盖局部损害。
\item 只在证据支持的范围内扩大使用；知识、模型、工具或流程变化后对受影响任务重新评估。
\end{enumerate}
\end{bookalgorithm}



% BEGIN GENERATED INTERVIEWS

% END GENERATED INTERVIEWS
\chapterreferences
\end{document}
